\SourceLeaf{15}
\[
-\frac{\mathrm{d}s'.\Phi r'.\cos MQ'}{r'r'}-\frac{\mathrm{d}s''.\Phi r''.\cos MQ''}{r''r''}-\frac{\mathrm{d}s'''.\Phi r'''.\cos MQ'''}{r'''r'''}-\text{ etc.}
\]
\par
and in the latter case by the same formula, adding the term
\[
-\mathrm{d}\Sigma.\Phi0
\]
\par
Multiplying this expression by $\cos MX$, we obtain the force with which the parts of the body inside the cone attract the point in the direction parallel and opposite to the $x$-coordinate axis. Hence the force with which the whole body acts in the same direction is expressed by the integral $-\int\frac{\mathrm{d}s.\Phi r.\cos MQ.\cos MX}{rr}$ taken over the whole surface of the body, provided the attracted point lies outside the body. But the integral $-\Phi0.\int\mathrm{d}\Sigma.\cos MX$ taken over the whole spherical surface must also be added whenever $M$ lies inside the body. It is also easy to see that, when $M$ lies on the surface of the body, the same integral $-\Phi0.\int\mathrm{d}\Sigma.\cos MX$ must indeed be added, but taken over only half the surface of the sphere. It is taken over the hemisphere bounded by the plane tangent to the surface of the body at $M$ and lying on the same side of that plane as the solid part of the body at point $M$. To determine the value of this integral, let us imagine the solid enclosed by that hemisphere and the plane. Let $\theta$ denote, in general, the angle between a line normal to the surface of this solid and directed outward and a line parallel to the $x$-coordinate axis. By Theorem One the integral $\int\mathrm{d}s.\cos\theta$ taken over the \textit{whole} surface of the solid vanishes. Thus, if the integral taken over the flat part of the surface alone is assumed to be $=J$, the integral taken over the curved part must be $=-J$. But on the curved part $\mathrm{d}s$ coincides with our $\mathrm{d}\Sigma$, and $\theta$ is $=180^0-MX$. Hence the integral $-\int\mathrm{d}\Sigma.\cos MX$ taken over the hemisphere is $=-J$. On the flat part of the surface, however, $\theta$ is evidently constant and equal to the value of $QX$ at point $M$. Thus $J$ is equal to the cosine of this angle multiplied by the area of the plane region, which is $=\pi$. It follows that the integral $-\Phi0.\int\mathrm{d}\Sigma.\cos MX$ taken over the hemisphere defined above is $=-\pi\Phi0.\cos QX$, with $QX$ evaluated at point $M$. In exactly the same way we find that the value of the integral $-\Phi0.\int\mathrm{d}\Sigma.\cos MX$ taken over the other hemisphere is $=+\pi\Phi0.\cos QX$, so that the integral over the whole sphere is $=0$. From all this we obtain
